\begin{minipage}{0.55\linewidth}
    On réalise une expérience en introduisant, à l'instant $t_{0}=0$, une masse de
    zinc en poudre de valeur $m(\ce{Zn})=\qty{1.0}{\g}$ dans un ballon contenant le
    volume $V=\qty{40}{\mL}$ d'une solution aqueuse ($S$) d'acide chlorhydrique
    (\ce{H3O^{+}_{(aq)} + Cl^{-}_{(aq)}}) de concentration molaire
    $C_{A}=\qty{0.5}{\mol\per\L}$. Les ions \ce{H3O^{+}_{(aq)}} réagissent avec le
    zinc \ce{Zn_{(s)}} suivant la réaction chimique d'équation~:
    \begin{center}
        \ce{2H3O^{+}_{(aq)} + Zn_{(s)} -> H2_{(g)} + Zn^{2+}_{(aq)} + 2H2O$_{(\ell)}$}.
    \end{center}
    
    La mesure du volume de dihydrogène formé permet le suivi de l'évolution
    temporelle de l'avancement $x$ de la réaction et de tracer le graphe $x=f(t)$.
    
    \begin{donnees}
        $M(\ce{Zn})=\qty{65.4}{\g\per\mol}$
    \end{donnees}
\end{minipage}
\hfill
\begin{minipage}{0.44\linewidth}
\begin{figure}[H]
    \centering
    \includegraphics[width=\textwidth]{2025-11-27_110541-}
\end{figure}
\end{minipage}

\begin{enumerate}
    \item Calculer les quantités de matière $n_{0}(Zn)$ et $n_{0}(\ce{H3O+})$,
          présentes initialement dans le mélange réactionnel.
    \item Dresser le tableau d'avancement de la réaction chimique.
    \item Identifier le réactif limitant. Justifier.
    \item Déterminer graphiquement~:
          \begin{enumerate}
              \item la valeur du temps de demi-réaction $t_{1/2}$.
              \item la valeur de la vitesse volumique de réaction, en unité
                    (\unit{\mol\per\L\per\s}), à l'instant $t=\qty{400}{\s}$,
                    sachant que le volume du mélange réactionnel est
                    $V=\qty{40}{\mL}$.
          \end{enumerate}
    \item Interpréter qualitativement la variation de la vitesse volumique de
          cette réaction.
    \item Pour accélérer la réaction précédente, on recommence l'expérience en
          utilisant la même masse de zinc $m(\ce{Zn})=\qty{1.0}{\g}$ et le
          volume $V=\qty{40}{\mL}$
          d'une solution aqueuse ($S^{\prime}$) d'acide chlorhydrique de
          concentration molaire $C_{A}{ }^{\prime}=\qty{1}{\mol\per\L}$.
          \begin{enumerate}
              \item Citer le facteur cinétique qui est à l'origine de
                    l'accélération de la réaction.
              \item Le temps de demi-réaction $t_{1/2}$ va-t-il augmenter ou
                    diminuer~? Justifier.
          \end{enumerate}
\end{enumerate}


